% Source: convolution_for_image_gradients.tex
\begin{notechapter}
\chapter{Image Gradients by Convolution}\label{ch:image-gradients}

\section{Introduction}
Convolution is a fundamental operation in image processing,
used to compute gradients for tasks like edge detection,
feature extraction, and optimization in machine learning.

\section{Gradient in an Image}
The gradient of an image represents the rate of change of
pixel intensity in different directions (e.g., horizontal
and vertical). It is a vector composed of partial derivatives
in the horizontal ($G_x$) and vertical ($G_y$) directions.

\section{Convolution for Gradient Computation}
Convolution applies filters (kernels) to an image to compute
gradients. A kernel is a small matrix that slides over the
image, performing element-wise multiplication and summation.

\subsection{Common Gradient Kernels}
\begin{itemize}
    \item \textbf{Sobel Operator}:
        \begin{align*}
            K_x &= \begin{bmatrix}
                    -1 & 0 & 1 \\
                    -2 & 0 & 2 \\
                    -1 & 0 & 1
                  \end{bmatrix}, \\
            K_y &= \begin{bmatrix}
                    -1 & -2 & -1 \\
                     0 &  0 &  0 \\
                     1 &  2 &  1
                  \end{bmatrix}.
        \end{align*}
    \item \textbf{Prewitt Operator}:
        \begin{align*}
            K_x &= \begin{bmatrix}
                    -1 & 0 & 1 \\
                    -1 & 0 & 1 \\
                    -1 & 0 & 1
                  \end{bmatrix}, \\
            K_y &= \begin{bmatrix}
                    -1 & -1 & -1 \\
                     0 &  0 &  0 \\
                     1 &  1 &  1
                  \end{bmatrix}.
        \end{align*}
    \item \textbf{Scharr Operator}:
        \begin{align*}
            K_x &= \begin{bmatrix}
                    -3 & 0 &  3 \\
                    -10 & 0 & 10 \\
                    -3 & 0 &  3
                  \end{bmatrix}, \\
            K_y &= \begin{bmatrix}
                    -3 & -10 & -3 \\
                     0 &   0 &  0 \\
                     3 &  10 &  3
                  \end{bmatrix}.
        \end{align*}
\end{itemize}

\section{Steps for Gradient Computation}
\begin{enumerate}
        \item \textbf{Convolve the Image}:
        Apply the horizontal kernel ($K_x$) to compute $G_x$,
        and the vertical kernel ($K_y$) to compute $G_y$.
        Signs depend on conventions: sliding \( K_x \) without flipping (cross-correlation, as in
        deep-learning libraries) computes ``right minus left'', i.e.\ a positive multiple of
        \( \partial I/\partial x \); a true convolution flips the kernel and gives the opposite sign.
        Likewise, with the image row index increasing downwards, \( K_y \) applied by cross-correlation
        computes ``bottom minus top''. The magnitude is unaffected by these sign conventions.
        \begin{itemize}
            \item The horizontal gradient $G_x$ detects vertical edges,
                  as it measures the rate of change in intensity along
                  the horizontal direction.
            \item The vertical gradient $G_y$ detects horizontal edges,
                  as it measures the rate of change in intensity along
                  the vertical direction.
            \item To ensure $G_x$ and $G_y$ have the same dimensions
                  as the input image $I$, \textbf{zero-padding} is used.
            \item Zero-padding adds a border of zeros around the image
                  to handle edge pixels during convolution.
        \end{itemize}
    \item \textbf{Combine Gradients}:
        Compute the gradient magnitude $G$ and direction $\theta$:
        \begin{align*}
            G &= \sqrt{G_x^2 + G_y^2}, \\
                        \theta &= \operatorname{atan2}(G_y, G_x).
        \end{align*}
        (The two-argument arctangent gives the full range \( (-180^\circ, 180^\circ] \); \( \arctan(G_y/G_x) \)
        only gives \( (-90^\circ, 90^\circ) \) and fails when \( G_x = 0 \).)
        \begin{itemize}
            \item The gradient magnitude $G$ represents the strength
                  of the intensity change at each pixel.
            \item The gradient direction $\theta$ represents the
                  orientation of the intensity change, which is
                  perpendicular to the edge.
        \end{itemize}
    \item \textbf{Visualize Gradients}:
        The magnitude image highlights edges, while the
        direction image shows the orientation of edges.
        \begin{itemize}
            \item The direction image is often visualized using a
                  colormap (e.g., HSV) to represent angles.
            \item With the usual linear mapping of the angle \( \theta \in [0^\circ, 360^\circ) \) to the hue
                  of the HSV color wheel, each color corresponds to an orientation:
                  \begin{itemize}
                      \item Red: $0^\circ$ (horizontal, rightward).
                      \item Yellow: $60^\circ$.
                      \item Green: $120^\circ$.
                      \item Cyan: $180^\circ$ (horizontal, leftward).
                      \item Blue: $240^\circ$.
                      \item Magenta: $300^\circ$.
                  \end{itemize}
                  Opposite directions (\( \theta \) and \( \theta + 180^\circ \)) receive complementary colors
                  (e.g.\ red and cyan).
        \end{itemize}
\end{enumerate}

\section{Feature Extraction Using Image Gradients}
Image gradients are widely used for feature extraction in
computer vision. They capture important information about
the structure and texture of an image, such as edges, corners,
and keypoints. These features are used for tasks like object
detection, image matching, and classification.

\subsection{Methods for Feature Extraction}
\begin{itemize}
    \item \textbf{Edge Detection}:
        Gradients are directly used to detect edges in an image.
        Algorithms like the \textbf{Canny edge detector} use
        gradient magnitude and direction to identify edges.

    \item \textbf{Corner Detection}:
        Corners are points where the intensity changes significantly
        in multiple directions. Gradient-based methods like the
        \textbf{Harris corner detector} use the gradient structure
        tensor to identify corners:
        \[
        M = \begin{bmatrix}
        \sum G_x^2 & \sum G_x G_y \\
        \sum G_x G_y & \sum G_y^2
        \end{bmatrix}.
        \]
        The eigenvalues of $M$ determine whether a point is a corner,
        edge, or flat region.

    \item \textbf{Histogram of Oriented Gradients (HOG)}:
        HOG is a popular feature descriptor that uses gradient
        orientation and magnitude to describe local shape information.
        Steps:
        \begin{enumerate}
            \item Compute gradients $G_x$ and $G_y$ for the image.
            \item Calculate gradient magnitude $G$ and direction $\theta$.
            \item Divide the image into small cells and create a histogram
                  of gradient orientations for each cell.
            \item Normalize the histograms across blocks of cells.
            \item Concatenate all histograms to form the final feature vector.
        \end{enumerate}
        HOG is widely used in object detection, such as pedestrian detection.

    \item \textbf{Scale-Invariant Feature Transform (SIFT)}:
        SIFT is a robust feature extraction algorithm that uses gradients
        to detect and describe keypoints. Steps:
        \begin{enumerate}
            \item Detect keypoints using the Difference of Gaussians (DoG) method.
            \item Compute gradients $G_x$ and $G_y$ in the region around each keypoint.
            \item Create a gradient orientation histogram for each keypoint.
            \item Use the histogram to generate a descriptor that is invariant
                  to rotation, scaling, and lighting changes.
        \end{enumerate}
        SIFT is widely used for image matching and object recognition.

    \item \textbf{Texture Analysis}:
        Gradients can be used to analyze texture by capturing patterns
        of intensity changes. For example, the \textbf{Local Binary Patterns (LBP)}
        method can be combined with gradient information to describe texture features.
\end{itemize}

\section{Applications of Image Gradients}
\begin{itemize}
    \item Edge detection
    \item Feature extraction (e.g., SIFT, HOG)
    \item Optimization in CNNs
    \item Image enhancement
\end{itemize}

\section{Mathematical Interpretation}
Convolution for gradient computation approximates the
derivative of the image intensity function $I(x, y)$ using
finite differences:
\begin{align*}
    G_x &= \frac{\partial I}{\partial x} \approx I * K_x, \\
    G_y &= \frac{\partial I}{\partial y} \approx I * K_y.
\end{align*}
Here, $*$ denotes the convolution operation, and $K_x$ and $K_y$
are finite difference kernels.

\section{Python Example}
Below is an example of computing gradient magnitude and direction
using OpenCV:
\begin{verbatim}
import cv2
import numpy as np
import matplotlib.pyplot as plt

# Load image in grayscale
image = cv2.imread('image.jpg', cv2.IMREAD_GRAYSCALE)

# Apply Sobel operator with zero-padding
G_x = cv2.Sobel(image, cv2.CV_64F, 1, 0, ksize=3)
G_y = cv2.Sobel(image, cv2.CV_64F, 0, 1, ksize=3)

# Compute gradient magnitude and direction
magnitude = np.sqrt(G_x**2 + G_y**2)
direction = np.arctan2(G_y, G_x)  # Direction in radians (-π to π)

# Normalize direction for visualization
direction_normalized = (direction + np.pi) * (255 / (2 * np.pi))
direction_normalized = direction_normalized.astype(np.uint8)

# Display the direction image
plt.imshow(direction_normalized, cmap='hsv')  # Use HSV colormap
plt.colorbar(label='Gradient Direction (radians)')
plt.title('Gradient Direction')
plt.show()
\end{verbatim}

\section{Conclusion}
Convolution is a powerful tool for computing image gradients,
enabling tasks like edge detection, feature extraction, and
optimization in machine learning models. By applying specific
kernels and using zero-padding, we can approximate the partial
derivatives of an image and extract meaningful information
about its structure.
\end{notechapter}
